\documentclass{amsart}
% For dealing with references we use the comment environment
\usepackage{verbatim}
\newenvironment{reference}{\comment}{\endcomment}
%\newenvironment{reference}{}{}
\newenvironment{slogan}{\comment}{\endcomment}
\newenvironment{history}{\comment}{\endcomment}

% For commutative diagrams we use Xy-pic
\usepackage[all]{xy}

% We use 2cell for 2-commutative diagrams.
\xyoption{2cell}
\UseAllTwocells

% We use multicol for the list of chapters between chapters
\usepackage{multicol}

% This is generally recommended for better output
\usepackage{lmodern}
\usepackage[T1]{fontenc}

% For cross-file-references

% Package for hypertext links:
\usepackage{hyperref}

% For any local file, say "hello.tex" you want to link to please

% Theorem environments.
%
\theoremstyle{plain}
\newtheorem{theorem}[subsection]{Theorem}
\newtheorem{proposition}[subsection]{Proposition}
\newtheorem{lemma}[subsection]{Lemma}

\theoremstyle{definition}
\newtheorem{definition}[subsection]{Definition}
\newtheorem{example}[subsection]{Example}
\newtheorem{exercise}[subsection]{Exercise}
\newtheorem{situation}[subsection]{Situation}

\theoremstyle{remark}
\newtheorem{remark}[subsection]{Remark}
\newtheorem{remarks}[subsection]{Remarks}

\numberwithin{equation}{subsection}

% Macros
%
\def\lim{\mathop{\mathrm{lim}}\nolimits}
\def\colim{\mathop{\mathrm{colim}}\nolimits}
\def\Spec{\mathop{\mathrm{Spec}}}
\def\Hom{\mathop{\mathrm{Hom}}\nolimits}
\def\Ext{\mathop{\mathrm{Ext}}\nolimits}
\def\SheafHom{\mathop{\mathcal{H}\!\mathit{om}}\nolimits}
\def\SheafExt{\mathop{\mathcal{E}\!\mathit{xt}}\nolimits}
\def\Sch{\mathit{Sch}}
\def\Mor{\mathop{\mathrm{Mor}}\nolimits}
\def\Ob{\mathop{\mathrm{Ob}}\nolimits}
\def\Sh{\mathop{\mathit{Sh}}\nolimits}
\def\NL{\mathop{N\!L}\nolimits}
\def\CH{\mathop{\mathrm{CH}}\nolimits}
\def\proetale{{pro\text{-}\acute{e}tale}}
\def\etale{{\acute{e}tale}}
\def\QCoh{\mathit{QCoh}}
\def\Ker{\mathop{\mathrm{Ker}}}
\def\Im{\mathop{\mathrm{Im}}}
\def\Coker{\mathop{\mathrm{Coker}}}
\def\Coim{\mathop{\mathrm{Coim}}}

% Boxtimes
%
\DeclareMathSymbol{\boxtimes}{\mathbin}{AMSa}{"02}

%
% Macros for moduli stacks/spaces
%
\def\QCohstack{\mathcal{QC}\!\mathit{oh}}
\def\Cohstack{\mathcal{C}\!\mathit{oh}}
\def\Spacesstack{\mathcal{S}\!\mathit{paces}}
\def\Quotfunctor{\mathrm{Quot}}
\def\Hilbfunctor{\mathrm{Hilb}}
\def\Curvesstack{\mathcal{C}\!\mathit{urves}}
\def\Polarizedstack{\mathcal{P}\!\mathit{olarized}}
\def\Complexesstack{\mathcal{C}\!\mathit{omplexes}}
% \Pic is the operator that assigns to X its picard group, usage \Pic(X)
% \Picardstack_{X/B} denotes the Picard stack of X over B
% \Picardfunctor_{X/B} denotes the Picard functor of X over B
\def\Pic{\mathop{\mathrm{Pic}}\nolimits}
\def\Picardstack{\mathcal{P}\!\mathit{ic}}
\def\Picardfunctor{\mathrm{Pic}}
\def\Deformationcategory{\mathcal{D}\!\mathit{ef}}

\setlength{\emergencystretch}{3em}
\begin{document}
\title[Stacks Project, Tag 0EXX]{422 --- Tame ramification is preserved under normal closure and compositum}
\maketitle
\noindent Copyright (C) 2005--2025 Johan de Jong. Stacks Project authors. Source: \href{https://stacks.math.columbia.edu/tag/0EXX}{Tag 0EXX}.

\noindent Editorial difficulty note (not author text). Difficulty H2: The proof reduces tame extensions to unramified extensions after adjoining a root of a uniformizer. For normal closure it explicitly analyzes normalized tensor powers, constructs the root-of-unity ratio algebra, and proves that all resulting field factors remain tame..

\noindent Editorial provenance note (not author text). History: excerpt from revision \texttt{a0444\allowbreak 6e57e\allowbreak c1fbc\allowbreak 252a8\allowbreak 71afc\allowbreak ec775\allowbreak 2fb28\allowbreak 07b14}, more-algebra.tex, lines 34242--34327. Reference presentation was adapted; original-author context and core support blocks were retained or added. The mathematical wording is unchanged. Licensed under GNU FDL 1.2 or later; no invariant sections or cover texts. See ../licenses/COPYING. Context blocks reproduce author definitions and supporting results; they are not additional counted problems.

\begin{definition}
\label{fields-definition-normal-closure}
Let $E/F$ be a finite extension of fields. The field extension $K/E$
constructed in Lemma \href{https://stacks.math.columbia.edu/tag/09DT}{lemma normal closure (Tag 09DT)}
is called the {\it normal closure $E$ over $F$}.
\end{definition}

\begin{lemma}
\label{lemma-permanence-tame}
Let $A$ be a discrete valuation ring with fraction field $K$.
\begin{enumerate}
\item If $L/K$ is a finite separable extension which is tamely
ramified with respect to $A$, then there exists a Galois
extension $M/K$ containing $L$ which is tamely ramified
with respect to $A$.
\item If $L_1/K$, $L_2/K$ are finite separable extensions which are tamely
ramified with respect to $A$, then there exists a finite
separable extension $L/K$ which is tamely ramified with respect
to $A$ containing $L_1$ and $L_2$.
\end{enumerate}
\end{lemma}

\begin{proof}
Proof of (2). Choose a uniformizer $\pi \in A$.
We can choose an integer $e$ invertible
in $\kappa_A$ and extensions $L_i'/K' = K[\pi^{1/e}]$
unramified with respect to $A' = A[\pi^{1/e}]$
with $L'_i/L_i$ as extensions of $K$, see
Lemma \href{https://stacks.math.columbia.edu/tag/0EXW}{lemma characterize tame (Tag 0EXW)}.
By Lemma \href{https://stacks.math.columbia.edu/tag/0EXR}{lemma permanence unramified (Tag 0EXR)}
we can find an extension $L'/K'$ which is unramified
with respect to $A'$ such that $L'_i/K$ is isomorphic
to a subextension of $L'/K'$ for $i = 1, 2$.
This finishes the proof of (3) as $L'/K$ is tamely ramified
(use same lemma as above).

\medskip\noindent
Proof of (1). We may first replace $L$ by a larger extension
and assume that $L$ is an extension of $K' = K[\pi^{1/e}]$
unramified with respect to $A' = A[\pi^{1/e}]$ where $e$ is
invertible in $\kappa_A$, see Lemma \href{https://stacks.math.columbia.edu/tag/0EXW}{lemma characterize tame (Tag 0EXW)}.
Let $M$ be the normal closure of $L$ over $K$, see
Fields, Definition \ref{fields-definition-normal-closure}.
Then $M/K$ is Galois by Fields, Lemma \href{https://stacks.math.columbia.edu/tag/0EXM}{lemma normal closure galois (Tag 0EXM)}.
On the other hand, there is a surjection
$$
L \otimes_K \ldots \otimes_K L \longrightarrow M
$$
of $K$-algebras, see Fields, Lemma
\href{https://stacks.math.columbia.edu/tag/0EXL}{lemma normal closure tensor product (Tag 0EXL)}.
Let $B$ be the integral closure of $A$ in $L$
as in Remark \href{https://stacks.math.columbia.edu/tag/09E8}{remark finite separable extension (Tag 09E8)}. The
condition that $L$ is unramified with respect to $A' = A[\pi^{1/e}]$
exactly means that $A' \to B$ is an \'etale ring map, see
Algebra, Lemma \href{https://stacks.math.columbia.edu/tag/00U6}{lemma characterize etale (Tag 00U6)}.
Claim: 
$$
K' \otimes_K \ldots \otimes_K K' = \prod K'_i
$$
is a product of field extensions $K'_i/K$ tamely
ramified with respect to $A$. Then if $A'_i$ is the integral
closure of $A$ in $K'_i$ we see that
$$
\prod A'_i \otimes_{(A' \otimes_A \ldots \otimes_A A')}
(B \otimes_A \ldots \otimes_A B)
$$
is finite \'etale over $\prod A'_i$ and hence a product of
Dedekind domains (Lemma \href{https://stacks.math.columbia.edu/tag/0AP2}{lemma Dedekind etale extension (Tag 0AP2)}).
We conclude that $M$ is the fraction field
of one of these Dedekind domains which is finite \'etale
over $A'_i$ for some $i$. It follows that $M/K'_i$
is unramified with respect to every maximal ideal of $A'_i$
and hence $M/K$ is tamely ramified by Lemma \href{https://stacks.math.columbia.edu/tag/0EXU}{lemma composition tame (Tag 0EXU)}.

\medskip\noindent
It remains the prove the claim. For this we write
$A' = A[x]/(x^e - \pi)$ and we see that
$$
A' \otimes_A \ldots \otimes_A A' =
A'[x_1, \ldots, x_r]/(x_1^e - \pi, \ldots, x_r^e - \pi)
$$
The normalization of this ring certainly contains the
elements $y_i = x_i/x_1$ for $i = 2, \ldots, r$ subject
to the relations $y_i^e - 1 = 0$ and
we obtain
$$
A[x_1, y_2, \ldots, y_r]/(x_1^e - \pi, y_2^e - 1, \ldots, y_r^e - 1) =
A'[y_2, \ldots, y_r]/(y_2^e - 1, \ldots, y_r^e - 1)
$$
This ring is finite \'etale over $A'$ because $e$ is invertible in $A'$.
Hence it is a product of Dedekind domains each unramified over $A'$
as desired (see references given above in case of confusion).
\end{proof}
\end{document}
